\chapter{Data Collection and the Riot API}

The system's primary data source is the Riot Developer Portal and the associated League of Legends API suite \cite{riotportal,riotapi}. These enable identifying player accounts, querying match lists, downloading detailed match JSON files, and processing participant statistics.

\section{Logic of the Riot API Endpoints Used}

The current project uses several layers of the Riot API, but the most important design decision was making \textit{puuid} the central identifier of the system. The pipeline breaks down logically into the following steps:

\begin{enumerate}
    \item starting from an initial player or seed list, the system requests account information;
    \item from the account, it retrieves a list of match identifiers;
    \item for the match identifiers it downloads the detailed match JSON files;
    \item from the participant rows it constructs a \textit{puuid}-based player profile and relationship data;
    \item the results are persisted bound to the dataset.
\end{enumerate}

The collector calls the \textit{match-v5} endpoint family via the \textit{europe.api.riotgames.com} regional route.

\section{Collector State Machine and Premade-Probe}

\textit{backend/match\_collector.py} is a parameterisable state machine: its configuration holds the dataset identifier, paths, collection mode, queue type, match count limit, iteration limit, and premade threshold.

\begin{figure}[H]
\centering
\includegraphics[width=0.98\textwidth,height=0.74\textheight,keepaspectratio]{collector-state-machine.pdf}
\caption{Flow diagram of the collector's parameterisable state machine}
\end{figure}

\section{Collection Focus and Queue Selection}

The project uses two research datasets: \textit{flexset} consists of Ranked Flex Queue matches, and \textit{soloq} consists of Ranked Solo/Duo Queue matches. The queue type choice is not arbitrary. Ranked Flex allows premade teams to participate, so recurring ally relationships can appear naturally. Ranked Solo/Duo intentionally excludes 5-player premades, making it a controlled comparison baseline.

\section{Rate Limits and Reliability}

The Riot API's constraints also shape the architecture. The collection layer gained explicit rate-limit handling, wait logic, and subtask splitting. The collector applies its own per-second wait, waits thirty seconds on HTTP 429 before retrying, and skips already-saved matches based on file existence.

\begin{table}[H]
\centering
\caption{Main reliability mechanisms of the collector}
\begin{tabularx}{\textwidth}{|>{\raggedright\arraybackslash}p{4.4cm}|X|}
\hline
\textbf{Mechanism} & \textbf{Function} \\
\hline
Per-second delay & Controls short-term API load \\
Two-minute counter & Aligns with Riot key usage limits \\
429 retry & More robust operation on transient limit hits \\
Existing match file skip & Avoids duplicate saves and re-downloads \\
Event-based stdout log & Supports status tracking on both backend and frontend \\
\hline
\end{tabularx}
\end{table}

These mechanisms are not mere technical details: collecting the full \textit{flexset} ($\approx$4981 matches) with a development API key took days, with interruptions. Without file-based skipping and restartability, the corpus could not have been assembled without a continuously running server. The collector's design decisions are therefore a direct response to the API constraints: iterated, stateful batch runs rather than a single-pass download. After each overnight run, collection resumed exactly where it left off.

\section{Comparison of flexset and soloq Collection Strategies}

The \textit{flexset} dataset contains exclusively queue ID 440 (Ranked Flex) matches from EUNE Master$+$ ladder seed players. The collector's premade-probe mechanism filters out candidates that lack sufficient recurring teammate patterns.

The \textit{soloq} dataset consists of queue ID 420 (Ranked Solo/Duo) matches from an EUNE Master$+$ Solo/Duo ladder seed population. The role of \textit{soloq} is to serve as a control baseline: a dataset in which recurring premade connections are systematically not favoured.

\begin{table}[H]
\centering
\caption{Comparison of the two dataset collection strategies}
\begin{tabularx}{\textwidth}{|>{\raggedright\arraybackslash}p{4.2cm}|X|X|}
\hline
\textbf{Dimension} & \textbf{flexset} & \textbf{soloq} \\
\hline
Queue ID & 440 (Flex) & 420 (Solo/Duo) \\
Premade-probe & Yes, based on repetitions & No \\
Seed expansion & Via teammates & Via list retrieval \\
Collection goal & Social structure & Performance baseline \\
Corpus size & 4980 matches & 2000 matches \\
Role in project & Main research dataset & Control dataset \\
\hline
\end{tabularx}
\end{table}

When interpreting Riot API data it is important to keep in mind that the network is \emph{not a representation of the entire League of Legends universe}, but a bounded sub-world created by a specific collection strategy. Detailed data quality and bias analysis can be found in chapter 4 of the extended version \cite{premadegraph-full}.
